\documentclass[11pt]{article}
\usepackage{amsmath,amssymb,amsthm}
\usepackage[margin=1in]{geometry}

\newtheorem{theorem}{Theorem}
\newtheorem{lemma}[theorem]{Lemma}

\title{Problem 432: Totient Stairstep Sequences}
\author{Project Euler Solutions}
\date{}

\begin{document}
\maketitle

\section{Problem Statement}
Define the iterated Euler totient function: $\phi^{(0)}(n) = n$, $\phi^{(k)}(n) = \phi(\phi^{(k-1)}(n))$. Let $f(n) = \min\{k : \phi^{(k)}(n) = 1\}$ with $f(1) = 0$. Find $\sum_{n=2}^{N} f(n)$ for $N = 5 \times 10^8$.

\section{Mathematical Foundation}

\begin{theorem}[Strict Decrease of Totient]
For every integer $n \geq 2$, $\phi(n) < n$.
\end{theorem}

\begin{proof}
Since $n \geq 2$, there exists a prime $p \mid n$. Among $\{1, \ldots, n\}$, the $n/p$ multiples of $p$ are not coprime to $n$, so $\phi(n) \leq n - n/p < n$.
\end{proof}

\begin{lemma}[Well-Definedness of $f$]
The function $f(n)$ is well-defined for all $n \geq 1$.
\end{lemma}

\begin{proof}
By Theorem~1, $\phi^{(k)}(n)$ is a strictly decreasing sequence of positive integers for $n \geq 2$. Every such sequence is finite, so $\phi^{(k)}(n) = 1$ for some $k$.
\end{proof}

\begin{theorem}[Recurrence]
For $n \geq 2$, $f(n) = 1 + f(\phi(n))$.
\end{theorem}

\begin{proof}
One application of $\phi$ advances the chain by one step. The remaining steps from $\phi(n)$ to $1$ number $f(\phi(n))$ by definition.
\end{proof}

\begin{theorem}[Euler's Product Formula]
For $n = p_1^{a_1} \cdots p_k^{a_k}$,
\[
\phi(n) = n \prod_{p \mid n} \left(1 - \frac{1}{p}\right).
\]
\end{theorem}

\begin{proof}
Follows from the Chinese Remainder Theorem and $\phi(p^a) = p^{a-1}(p-1)$.
\end{proof}

\begin{lemma}[Sieve Correctness]
Initializing $\phi(n) = n$ and, for each prime $p$, multiplying $\phi(kp)$ by $(1 - 1/p)$ for all multiples $kp \leq N$ yields the correct totient for every $n \leq N$.
\end{lemma}

\begin{proof}
Each prime factor $p$ of $n$ contributes exactly one factor of $(1 - 1/p)$ to the product formula. The sieve visits each prime and applies this factor to all its multiples, producing the correct result.
\end{proof}

\section{Algorithm}
\begin{enumerate}
    \item Compute $\phi(n)$ for all $n \leq N$ via the Euler sieve.
    \item Compute $f(n) = 1 + f(\phi(n))$ for $n = 2, \ldots, N$ with $f(1) = 0$.
    \item Return $\sum_{n=2}^{N} f(n)$.
\end{enumerate}

\section{Complexity Analysis}
\begin{itemize}
    \item \textbf{Time:} $O(N \log\log N)$ for the sieve; $O(N)$ for computing $f$ and summation. Total: $O(N \log\log N)$.
    \item \textbf{Space:} $O(N)$ for arrays $\phi$ and $f$.
\end{itemize}

\section{Answer}

\[
\boxed{754862080}
\]

\end{document}
